\label{sec:tree-wave}
\subsection{Finite radial reduction}

Let $G_L$ be the depth-$L$ rooted $3$-regular tree.  Its levels
$\mathcal L_k$ have
\[
 |\mathcal L_0|=1,
 \qquad
 |\mathcal L_k|=3\cdot2^{k-1}\quad(k\ge1).
\]
The root and nonleaf vertices have degree three and leaves have degree one.
Use the root as the single source.  Radial symmetry implies that every
proximal iterate is constant on each level.  Define its total
degree-weighted shell mass and shell volume by
\begin{equation}
 m_{t,k}:=\sum_{v\in\mathcal L_k}\sqrt{d_v}\,x_{t,v},
 \qquad
 v_k:=\vol(\mathcal L_k).
 \label{eq:shell-mass}
\end{equation}
Then $p_t=\sum_{k=0}^Lm_{t,k}$.

Let $P_L$ be the column-stochastic radial random-walk matrix:
\begin{align*}
 (P_L)_{1,0}&=1,\\
 (P_L)_{k-1,k}&=\frac13,
 & (P_L)_{k+1,k}&=\frac23
 &&(1\le k<L),\\
 (P_L)_{L-1,L}&=1.
\end{align*}
For $\kappa=1-2\alpha$,
\begin{equation}
 Q+\kappa I
 =\frac{1-\alpha}{2}(3I-W).
 \label{eq:tree-prox-hessian}
\end{equation}
The similarity transformation from radial coordinate values to shell masses
turns $W$ into $P_L$.

\begin{proposition}[Exact radial obstacle recurrence; \ProvedStatus]
\label{prop:radial-recurrence}
The exact proximal iterate~\eqref{eq:app-step} on $G_L$ is the unique fixed
point of
\begin{equation}
 m_t=pos{
   \frac13P_Lm_t
   +\frac{2}{3(1-\alpha)}
    \left[
      \alpha e_0-\alpha\bar\rho v
      +\kappa\bigl((1+\beta)m_{t-1}-\beta m_{t-2}\bigr)
    \right]
 }.
 \label{eq:radial-obstacle}
\end{equation}
The positive part is componentwise.  In shell variables,
\begin{equation}
 d_t
 =\sum_{\substack{k:m_{t-2,k}>0\\m_{t,k}=0}}m_{t-2,k}.
 \label{eq:radial-dropout}
\end{equation}
\end{proposition}

\begin{proof}
Multiply the radial KKT system by the shell mass weights.  The quadratic term
is~\eqref{eq:tree-prox-hessian}, the source has mass $\alpha e_0$, and the
$\ell_1$ derivative has shell mass $\alpha\bar\rho v$.  Solving the obstacle
equation for $m_t$ gives~\eqref{eq:radial-obstacle}.  Since
$\norm{P_L}_1=1$ and positive-part projection is nonexpansive, its right side
is a $1/3$ contraction in $\ell_1$.  The fixed point is therefore unique.
Equation~\eqref{eq:radial-dropout} is~\eqref{eq:dropout-mass} grouped by
levels.
\end{proof}

This recurrence removes the main numerical ambiguity of the full graph: a
depth-$L$ tree with exponentially many vertices is represented by only $L+1$
nonnegative masses, and every proximal subproblem is solved by a stable
$1/3$-contractive iteration.

\subsection{Finite-tree stress test}

Fix
\begin{equation}
 \alpha=10^{-3},
 \qquad
 \bar\rho_L=10^{-0.275L}.
 \label{eq:tree-parameters}
\end{equation}
Since $1/\bar\rho_L=10^{0.275L}$ whereas
$|V(G_L)|=\Theta(2^L)=10^{0.3010\ldots L}$, this remains a genuinely local
parameter regime.

Starting from $m_{-1}=m_0=0$, we solved~\eqref{eq:radial-obstacle} until
successive fixed-point iterates differed by at most
$10^{-14}\max\{1,\norm{m_t}_1\}$.  The following values are direct
evaluations of~\eqref{eq:flux} and~\eqref{eq:radial-dropout}.

\begin{table}[ht]
\centering
\caption{Finite rooted-tree deactivation flux for~\eqref{eq:tree-parameters}.
These are computational observations, not hidden asymptotic claims.}
\label{tab:tree-flux}
\begin{tabular}{rrrr}
\toprule
Depth $L$ & Iteration $T$ & $\Phi_T$ & $\Phi_T/(\alpha T)$ \\
\midrule
60  & 117 & $3.552001$             & $3.0359\times10^{1}$ \\
100 & 198 & $6.291300\times10^{1}$ & $3.1774\times10^{2}$ \\
140 & 254 & $1.013375\times10^{3}$ & $3.9897\times10^{3}$ \\
180 & 335 & $1.602154\times10^{4}$ & $4.7826\times10^{4}$ \\
240 & 418 & $9.371581\times10^{5}$ & $2.2420\times10^{6}$ \\
280 & 498 & $1.392457\times10^{7}$ & $2.7961\times10^{7}$ \\
\bottomrule
\end{tabular}
\end{table}

At $L=100,T=198$,
\[
 \sum_{t\le198}d_t=67.024028\ldots,
 \qquad
 \Phi_{198}=62.913002\ldots,
 \qquad
 \alpha T=0.198.
\]
Changing the shell-support threshold between zero and $10^{-10}$ changes none
of these displayed digits.  Direct nonnegative quadratic programming in the
radial value coordinates agrees with the mass fixed point on moderate depths.

\paragraph{Interpretation.}
For each fixed finite graph, accelerated proximal point still converges in
objective value.  The instability is in a different norm: momentum launches
an outward radial wave whose Euclidean energy can remain controlled while its
degree-weighted $\ell_1$ mass grows across exponentially expanding shells.
Repeated shell disappearance then accumulates in $\Phi_T$.

\subsection{Exact outward traveling wave}

The mechanism can be isolated analytically.  On a rooted $B$-regular tree,
the inward current-shell coupling vanishes as $B\to\infty$.  Away from the
root and leaves,~\eqref{eq:radial-obstacle} becomes
\begin{equation}
 m_{t,k}
 =\pos{
   \frac13m_{t,k-1}
   +a_\alpha\bigl((1+\beta)m_{t-1,k}-\beta m_{t-2,k}\bigr)
 },
 \qquad
 a_\alpha:=\frac{2(1-2\alpha)}{3(1-\alpha)}.
 \label{eq:outward-recurrence}
\end{equation}
At $\alpha=0$ this is
\begin{equation}
 m_{t,k}
 =\pos{
   \frac13m_{t,k-1}
   +\frac43m_{t-1,k}-\frac23m_{t-2,k}
 }.
 \label{eq:zero-alpha-outward}
\end{equation}

\begin{proposition}[Growing period-11 wave; \ProvedStatus]
\label{prop:period-eleven-wave}
Let $g=1.186095084\ldots$ be the largest real root with the sign pattern
below of
\begin{equation}
 19683g^6-69984g^5+91854g^4-54000g^3
 +13560g^2-1088g+8=0.
 \label{eq:g-polynomial}
\end{equation}
Set $u_{-2}=u_{-1}=0$, $u_0=1$, and for $0\le j\le8$ define
\begin{equation}
 u_{j+1}=3g u_j-4g u_{j-1}+2g u_{j-2}.
 \label{eq:phase-recurrence}
\end{equation}
Then $u_9=u_{10}=0$, $u_0,\ldots,u_8>0$, and
\begin{equation}
 m_{t,k}=c g^k u_{(t-k)\bmod 11}
 \label{eq:traveling-wave}
\end{equation}
is an exact nonnegative solution of~\eqref{eq:zero-alpha-outward} for every
$c>0$.
\end{proposition}

\begin{proof}
The positive phases generated by~\eqref{eq:phase-recurrence} are
approximately
\[
 (u_0,\ldots,u_8)
 =(1,3.558,7.917,13.661,19.490,23.319,22.912,17.130,7.565).
\]
The equation $u_9=0$ factors as $g^3$ times
~\eqref{eq:g-polynomial}.  The two inactive phases are strict because
\[
 \frac43u_8-\frac23u_7<0,
 \qquad
 \frac{u_0}{3g}-\frac23u_8<0.
\]
Substitute~\eqref{eq:traveling-wave} into
~\eqref{eq:zero-alpha-outward}.  On phases $0,\ldots,8$, division by
$cg^{k-1}$ gives exactly~\eqref{eq:phase-recurrence}; on phases $9$ and $10$,
the two strict inequalities make the positive part zero.  The modulo wrap at
phase zero uses $u_1=3g u_0$.
\end{proof}

A shell positive at phase seven is zero two steps later at phase nine, so it
contributes mass proportional to $g^k u_7$ to a dropout.  The strict sign
pattern persists for small positive $\alpha$.  At $\alpha=10^{-3}$, the
analogous outward recurrence has a growth root
\begin{equation}
 1.1485<g_\alpha<1.1486.
 \label{eq:positive-alpha-growth}
\end{equation}

\subsection{Rigorous scope of the tree conclusion}

Three statements should not be conflated.
\begin{enumerate}[label=(\roman*),leftmargin=*]
  \item \cref{prop:radial-recurrence} is an exact finite-graph reduction.
  \item \cref{tab:tree-flux} is a reproducible numerical evaluation of that
  exact reduction and exhibits growth by more than six orders of magnitude in
  the normalized ratio.
  \item \cref{prop:period-eleven-wave} is an exact analytic growing solution
  of the outward limiting dynamics, and its strict phases persist under small
  coefficient perturbations.
\end{enumerate}

Together these facts invalidate~\eqref{eq:flux-conjecture} as a defensible
graph-uniform lemma for an unconditional algorithmic theorem: classical
inertia has a concrete branching-wave mechanism that the conjecture does not
charge.  A completely formal asymptotic disproof for the specific
\emph{zero-start, root-source, finite-tree} family still requires one final
transfer statement: prove that the causal root response has a uniformly
nonzero component in the growing phase and that finite branching and the leaf
boundary preserve it over the required horizon.  Alternatively, one can
produce interval-arithmetic certificates for an unbounded sequence of finite
depths.  The finite computations strongly support this transfer but are not
silently substituted for it here.

